\subtitle{\bf Aula 12 - Introdução}
\begin{frame}[t,plain]
\titlepage
\end{frame}

\begin{frame}{Objetivos}
\protect\hypertarget{objetivos-1}{}
\begin{itemize}
\item
  Apresentar os conceitos de semântica small e big step.
\item
  Mostrar como implementar um interpretador a partir da especificação da
  semântica.
\end{itemize}
\end{frame}

\hypertarget{introduuxe7uxe3o-uxe0-semuxe2ntica-formal}{%
\section{Introdução à Semântica
Formal}\label{introduuxe7uxe3o-uxe0-semuxe2ntica-formal}}

\begin{frame}{Semântica Formal}
\protect\hypertarget{semuxe2ntica-formal}{}
\begin{itemize}

\item
  Apos a análise léxica e sintática, sabemos se o programa está escrito
  corretamente.
\item
  A semântica responde: \textbf{o que o programa faz?}
\end{itemize}
\end{frame}

\begin{frame}{Semântica Formal}
\protect\hypertarget{semuxe2ntica-formal-1}{}
\begin{itemize}

\item
  \textbf{Semântica operacional}: define o significado de um programa
  descrevendo como ele é \emph{executado}
\item
  Especificação se traduz diretamente em código (interpretadores e
  compiladores)
\end{itemize}
\end{frame}

\begin{frame}{Dois Estilos}
\protect\hypertarget{dois-estilos}{}
\begin{itemize}

\item
  \textbf{Small-step} (\emph{pequenos passos}):

  \begin{itemize}

  \item
    Define cada passo elementar de redução
  \item
    Relação \(e \to e'\): ``a expressão \(e\) reduz em um passo para
    \(e'\)''
  \item
    Útil para raciocinar sobre propriedades de execução
  \end{itemize}
\end{itemize}
\end{frame}

\begin{frame}{Dois Estilos}
\protect\hypertarget{dois-estilos-1}{}
\begin{itemize}

\item
  \textbf{Big-step} (\emph{grandes passos}, ou semântica natural):

  \begin{itemize}

  \item
    Define diretamente o valor final
  \item
    Relação \(e \Downarrow v\): ``a expressão \(e\) avalia para o valor
    \(v\)''
  \item
    Mais direta para implementar interpretadores
  \end{itemize}
\end{itemize}
\end{frame}

\begin{frame}[fragile]{Expressões}
\protect\hypertarget{expressuxf5es}{}
\[\begin{array}{lcl}
e &::=& n \mid e_1 + e_2 \mid e_1 * e_2
\end{array}\]

Em Haskell:

\begin{lstlisting}[language=Haskell]
data Exp
  = EInt Int
  | Exp :+: Exp
  | Exp :*: Exp
\end{lstlisting}
\end{frame}

\hypertarget{semuxe2ntica-small-step}{%
\section{Semântica Small-Step}\label{semuxe2ntica-small-step}}

\begin{frame}{Valores}
\protect\hypertarget{valores}{}
\begin{itemize}

\item
  Um \textbf{valor} é uma expressão na forma final

  \begin{itemize}

  \item
    Para expressões aritméticas, é um número inteiro \(v ::= n\).
  \end{itemize}
\end{itemize}
\end{frame}

\begin{frame}{Regras}
\protect\hypertarget{regras}{}
\begin{itemize}

\item
  As regras de redução são apresentadas como sistemas de inferência:
\end{itemize}

\[\frac{\text{premissas}}{\text{conclusão}}\]
\end{frame}

\begin{frame}{Regras para Adição}
\protect\hypertarget{regras-para-adiuxe7uxe3o}{}
\begin{itemize}

\item
  \textbf{Adição de valores} (axioma):
\end{itemize}

\[\frac{}{n_1 + n_2 \to n_1 \mathbin{\hat{+}} n_2} \quad \text{(Add)}\]
\end{frame}

\begin{frame}{Regras para Adição}
\protect\hypertarget{regras-para-adiuxe7uxe3o-1}{}
\begin{itemize}

\item
  \textbf{Redução no lado esquerdo} (esquerda primeiro):
\end{itemize}

\[\frac{e_1 \to e_1'}{e_1 + e_2 \to e_1' + e_2} \quad \text{(AddL)}\]
\end{frame}

\begin{frame}{Regras para Adição}
\protect\hypertarget{regras-para-adiuxe7uxe3o-2}{}
\begin{itemize}

\item
  \textbf{Redução no lado direito} (quando esquerdo já é valor):
\end{itemize}

\[\frac{e_2 \to e_2'}{n_1 + e_2 \to n_1 + e_2'} \quad \text{(AddR)}\]
\end{frame}

\begin{frame}{Regras para Adição}
\protect\hypertarget{regras-para-adiuxe7uxe3o-3}{}
\begin{itemize}

\item
  A regra (AddR) só se aplica quando \(n_1\) é um valor

  \begin{itemize}

  \item
    Garante avaliação \textbf{da esquerda para a direita}.
  \end{itemize}
\end{itemize}
\end{frame}

\begin{frame}{Multiplicação}
\protect\hypertarget{multiplicauxe7uxe3o}{}
\[\frac{}{n_1 * n_2 \to n_1 \mathbin{\hat{*}} n_2} \quad \text{(Mul)}\]
\end{frame}

\begin{frame}{Multiplicação}
\protect\hypertarget{multiplicauxe7uxe3o-1}{}
\[\frac{e_1 \to e_1'}{e_1 * e_2 \to e_1' * e_2} \quad \text{(MulL)}\]
\end{frame}

\begin{frame}{Multiplicação}
\protect\hypertarget{multiplicauxe7uxe3o-2}{}
\[\frac{e_2 \to e_2'}{n_1 * e_2 \to n_1 * e_2'} \quad \text{(MulR)}\]
\end{frame}

\begin{frame}{Sequências de Redução}
\protect\hypertarget{sequuxeancias-de-reduuxe7uxe3o}{}
\begin{itemize}

\item
  A execução completa é uma sequência de passos terminando em um valor:
\end{itemize}

\[e \to e_1 \to e_2 \to \cdots \to v\]

\begin{itemize}

\item
  Escrevemos \(e \to^* v\) quando existe tal sequência (zero ou mais
  passos).
\end{itemize}
\end{frame}

\begin{frame}{Exemplo: Small-Step para \((2 + 3) * 4\)}
\protect\hypertarget{exemplo-small-step-para-2-3-4}{}
\[\begin{array}{ll}
(2 + 3) * 4 & \xrightarrow{\text{MulL (com Add)}} \\
5 * 4       & \xrightarrow{\text{Mul}} \\
20          &
\end{array}\]
\end{frame}

\begin{frame}{Exemplo: Small-Step para \(1 + 2 * 3\)}
\protect\hypertarget{exemplo-small-step-para-1-2-3}{}
\[\begin{array}{ll}
1 + 2 * 3 & \xrightarrow{\text{AddR (com Mul)}} \\
1 + 6     & \xrightarrow{\text{Add}} \\
7         &
\end{array}\]
\end{frame}

\hypertarget{semuxe2ntica-big-step}{%
\section{Semântica Big-Step}\label{semuxe2ntica-big-step}}

\begin{frame}{Regras Big-Step}
\protect\hypertarget{regras-big-step}{}
\begin{itemize}

\item
  A relação \(e \Downarrow v\) define diretamente o valor final.
\end{itemize}
\end{frame}

\begin{frame}{Regras Big-Step}
\protect\hypertarget{regras-big-step-1}{}
\begin{itemize}

\item
  \textbf{Um número avalia para si mesmo:}
\end{itemize}

\[\frac{}{n \Downarrow n} \quad \text{(Num)}\]
\end{frame}

\begin{frame}{Regras Big-Step}
\protect\hypertarget{regras-big-step-2}{}
\begin{itemize}

\item
  \textbf{Adição:}
\end{itemize}

\[\frac{e_1 \Downarrow n_1 \quad e_2 \Downarrow n_2}{e_1 + e_2 \Downarrow n_1 \mathbin{\hat{+}} n_2} \quad \text{(Add)}\]
\end{frame}

\begin{frame}{Regras Big-Step}
\protect\hypertarget{regras-big-step-3}{}
\begin{itemize}

\item
  \textbf{Multiplicação:}
\end{itemize}

\[\frac{e_1 \Downarrow n_1 \quad e_2 \Downarrow n_2}{e_1 * e_2 \Downarrow n_1 \mathbin{\hat{*}} n_2} \quad \text{(Mul)}\]
\end{frame}

\begin{frame}{Exemplo: Derivação Big-Step para \((2 + 3) * 4\)}
\protect\hypertarget{exemplo-derivauxe7uxe3o-big-step-para-2-3-4}{}
\[\dfrac{
  \dfrac{
    \dfrac{}{2 \Downarrow 2} \quad \dfrac{}{3 \Downarrow 3}
  }{2 + 3 \Downarrow 5}
  \quad
  \dfrac{}{4 \Downarrow 4}
}{(2 + 3) * 4 \Downarrow 20}\]
\end{frame}

\begin{frame}{Exemplo: Derivação Big-Step para \(1 + 2 * 3\)}
\protect\hypertarget{exemplo-derivauxe7uxe3o-big-step-para-1-2-3}{}
\[\dfrac{
  \dfrac{}{1 \Downarrow 1}
  \quad
  \dfrac{
    \dfrac{}{2 \Downarrow 2} \quad \dfrac{}{3 \Downarrow 3}
  }{2 * 3 \Downarrow 6}
}{1 + 2 * 3 \Downarrow 7}\]
\end{frame}

\hypertarget{implementauxe7uxe3o-em-haskell}{%
\section{Implementação em
Haskell}\label{implementauxe7uxe3o-em-haskell}}

\begin{frame}[fragile]{Interpretador}
\protect\hypertarget{interpretador}{}
\begin{itemize}

\item
  A semântica big-step se traduz em Haskell de forma quase literal:
\end{itemize}

\begin{lstlisting}[language=Haskell]
eval :: Exp -> Int
eval (EInt n)    = n
eval (e1 :+: e2) = eval e1 + eval e2
eval (e1 :*: e2) = eval e1 * eval e2
\end{lstlisting}
\end{frame}

\begin{frame}[fragile]{Interpretador}
\protect\hypertarget{interpretador-1}{}
\begin{itemize}
\item
  Correspondência com as regras:
\item
  \texttt{eval (EInt n) = n} implementa a regra
  \textbf{(Num)}
\item
  \texttt{eval (e1 :+: e2) = eval e1 + eval e2}
  implementa \textbf{(Add)}
\item
  \texttt{eval (e1 :*: e2) = eval e1 * eval e2}
  implementa \textbf{(Mul)}
\end{itemize}
\end{frame}

\begin{frame}[fragile]{Executando o Interpretador}
\protect\hypertarget{executando-o-interpretador}{}
\begin{lstlisting}[language=bash]
cabal run exp -- --interp --file input.exp
\end{lstlisting}

Para a entrada \texttt{(2 + 3) * 4}, o interpretador
produz \texttt{20}.

\begin{lstlisting}[language=Haskell]
-- O pipeline completo:
-- 1. Lexer: String -> [Token]
-- 2. Parser: [Token] -> Exp
-- 3. eval: Exp -> Int
\end{lstlisting}
\end{frame}

\hypertarget{conclusuxe3o}{%
\section{Conclusão}\label{conclusuxe3o}}

\begin{frame}{Conclusão}
\protect\hypertarget{conclusuxe3o-1}{}
\begin{itemize}

\item
  Apresentamos dois estilos de semântica operacional

  \begin{itemize}

  \item
    Big-step e small-step
  \end{itemize}
\item
  Mostramos como implementar um interpretador a partir de sua semântica
  big-step.
\end{itemize}
\end{frame}
